\chapter{Máquina de estados}

Este capítulo explica los estados por los que pasa una sesión y lo que se guarda en cada uno. La definición completa está en \texttt{docs/contexto/maquina-estados.md}. Aquí solo se describe lo implementado: la apertura (CU-04) y los turnos (CU-01 y CU-03).

\section{Estados de la sesión}

\begin{verbatim}
        OPEN_ROUND aceptado
               |
               v
            CREATED -- AgentCard no válida --> NOT_STARTED
               |                                 (final)
               | AgentCards válidas
               v
     +--> NEGOTIATING -- ACCEPT --> HELD
     |       |    |
     +-------+    | WITHDRAW, COUNTEROFFER del último
   PROPOSE o      | turno o primer fallo
   COUNTEROFFER   v
   (turn + 1)  CLOSING
                  | SESSION_CLOSE enviado o fallido
                  v
               CLOSED (final)
\end{verbatim}

\begin{tabular}{p{3cm}p{10.5cm}}
  \textbf{Estado} & \textbf{Significado} \\
  \texttt{CREATED} & Sesión creada por un \texttt{OPEN\_ROUND}; faltan por comprobar las AgentCards. \\
  \texttt{NOT\_STARTED} & Alguna AgentCard no se pudo obtener o no es válida (D-99, D-131). Es final. \\
  \texttt{NEGOTIATING} & Hay un turno en curso, en \texttt{READY} o en \texttt{AWAITING\_TASK}. \\
  \texttt{HELD} & Un \texttt{ACCEPT} dejó un acuerdo provisional sobre la oferta vigente (D-77, D-84). \\
  \texttt{CLOSING} & El resultado está decidido y guardado. Falta avisar a las partes con \texttt{SESSION\_CLOSE}. \\
  \texttt{CLOSED} & \texttt{SESSION\_CLOSE} se ha enviado o intentado a las dos partes. Es final, y el resultado no cambia (D-49). \\
\end{tabular}

La salida de \texttt{HELD} y el envío de \texttt{AGREEMENT\_HELD} están pendientes de implementar. Por ahora, una sesión en \texttt{HELD} se queda ahí, con ese envío marcado como \texttt{PENDIENTE}.

\section{Subestados del turno}

\begin{tabular}{p{3cm}p{10.5cm}}
  \textbf{Subestado} & \textbf{Qué hay guardado} \\
  \texttt{READY} & \texttt{turn}, \texttt{attempt} y el \texttt{messageId} del intento (D-66). Aún no se ha llamado al agente. \\
  \texttt{AWAITING\_TASK} & Además, el \texttt{taskId} que se sondea, el instante de inicio del intento y el \texttt{contextId} de la parte (D-58). \\
\end{tabular}

Solo se guardan los puntos en los que el orquestador espera a la red. La validación del acto ocurre en memoria, y su resultado se escribe directamente como la transición siguiente.

\section{Transiciones de un turno}

\begin{tabular}{p{5cm}p{8.5cm}}
  \textbf{Evento} & \textbf{Qué se escribe} \\
  \texttt{PROPOSE} o \texttt{COUNTEROFFER} con \texttt{turn} < \texttt{maxTurns} & Acto aceptado; \texttt{turn + 1} en \texttt{READY}, \texttt{messageId} nuevo y la oferta como oferta vigente, con su emisor. \\
  \texttt{COUNTEROFFER} en el último turno & Acto aceptado; \texttt{CLOSING} con \texttt{NO\_AGREEMENT} y \texttt{lastOffers}. La oferta vigente no cambia (D-54, D-57). \\
  \texttt{ACCEPT} & Acto aceptado; \texttt{HELD} con la oferta vigente como oferta acordada y \texttt{heldUntil} = instante del acto + \texttt{holdTimeoutSeconds} (D-86). \\
  \texttt{WITHDRAW} & Acto aceptado; \texttt{CLOSING} con \texttt{WITHDRAWN}, \texttt{withdrawnBy}, \texttt{reason} si la hay y \texttt{lastOffers} (D-82). \\
  Fallo del intento & Cada fallo al registro; \texttt{CLOSING} con \texttt{ABORTED}, \texttt{abortReason = RETRIES\_EXHAUSTED}, \texttt{lastFailure} y \texttt{lastOffers} (D-07, D-53, D-61). \\
\end{tabular}

\texttt{lastOffers} es la última oferta aceptada de cada parte que llegó a ofertar, sacada del registro de actos aceptados. Incluye el \texttt{COUNTEROFFER} del último turno. Si nadie ofertó, no aparece (D-82).

En esta versión, \texttt{maxRetries} es siempre 0, así que el primer fallo agota los reintentos y aborta la sesión. El reintento con \texttt{attempt + 1} y \texttt{retryReason} está pendiente de implementar (\texttt{feature/reintentos-aborto}).

\section{Cierre}

En \texttt{CLOSING}, el registro de sesión guarda el estado del envío de \texttt{SESSION\_CLOSE} a cada parte: \texttt{PENDIENTE}, \texttt{ENVIADO} o \texttt{FALLIDO}. Cada marca se escribe en su propia transacción, junto con su entrada en el registro de envíos de control (D-70). La transacción que marca el último envío pendiente pasa la sesión a \texttt{CLOSED} (CU-02). Si los dos envíos fallan, la sesión también llega a \texttt{CLOSED}, con el mismo resultado.

\section{Tablas}

\begin{tabular}{p{3.5cm}p{10cm}}
  \textbf{Tabla} & \textbf{Contenido} \\
  \texttt{sesiones} & Un registro por sesión, que se reescribe en cada transición. \\
  \texttt{actos\_aceptados} & Un acto por turno resuelto: \texttt{turn}, emisor, tipo, oferta o \texttt{reason}, \texttt{taskId} e instante. Solo se añade. \\
  \texttt{fallos} & Un fallo por fila: \texttt{turn}, \texttt{attempt}, \texttt{taskId} si lo hay, \texttt{code}, \texttt{field} e instante. Solo se añade. \\
  \texttt{coerciones} & Creada y vacía: solo se llena en modo tolerante, que aún no está implementado. \\
  \texttt{envios\_control} & Un envío de \texttt{SESSION\_CLOSE} por fila: parte, \texttt{turn}, instante, resultado y, según el caso, \texttt{taskId} y estado de la Task, si llegó un Message, o el código del fallo (D-70). Solo se añade. \\
\end{tabular}

Cada transición escribe el registro de sesión y sus filas de registro en una sola transacción (D-63, D-69). Todos los instantes los da el reloj de PostgreSQL.

\section{Concurrencia}

Una instancia solo trabaja con una sesión si tiene su arrendamiento (D-63, D-65). Lo renueva al tomar la sesión, en cada sondeo y en cada transición. Cada escritura exige que la versión del registro y la instancia sean las que leyó. Si otra instancia ha escrito antes, la escritura no afecta a ninguna fila y la instancia abandona la sesión sin escribir nada más.
